\chapter{DPF Literature-to-Debugging Crosswalk}
\label{ch:bf-dpf-debug-crosswalk}

The preceding DPF chapters now separate the main mathematical layers: particle
flow, proposal correction, differentiable resampling, OT/Sinkhorn relaxation,
learned transport, and HMC target interpretation.  This chapter turns that
separation into a debugging crosswalk.  Its purpose is deliberately practical:
when an implementation fails, the reader should be able to identify the likely
mathematical layer, return to the right chapter section, and select the next
diagnostic before changing code.

The crosswalk is not a substitute for unit tests or numerical experiments.  It
is a way to keep those tests aligned with the literature.  Particle degeneracy
should not be debugged as a Sinkhorn problem; a value-gradient inconsistency
should not be explained away as ordinary Monte Carlo noise; and a learned-OT
extrapolation failure should not be treated as evidence that the OT teacher map
itself is invalid.

\section{How to read a DPF failure}
\label{sec:bf-dpf-debug-how-to-read}

Most DPF failures are ambiguous at first observation.  A collapsed cloud, for
example, may arise from ordinary particle degeneracy, an overly aggressive flow,
soft-resampling shrinkage, an unstable learned map, or a target mismatch in the
HMC wrapper.  The first diagnostic step is therefore to locate the failure in
the computation graph:
\begin{enumerate}[label=(\roman*)]
  \item before the flow map;
  \item during the flow ODE;
  \item in proposal correction and log-determinant accounting;
  \item at the resampling or OT-relaxation layer;
  \item inside a learned transport approximation;
  \item or at the HMC value-gradient interface.
\end{enumerate}
Only after that localization should the implementation be changed.

\section{Failure-mode crosswalk}
\label{sec:bf-dpf-debug-crosswalk-table}

\begin{center}
\scriptsize
\begin{longtable}{@{}p{0.15\textwidth} p{0.17\textwidth} p{0.20\textwidth} p{0.18\textwidth} p{0.22\textwidth}@{}}
\caption{Literature-to-debugging crosswalk for the rebuilt DPF block.}
\label{tab:bf-dpf-debug-crosswalk}\\
\toprule
Failure mode & Likely mathematical layer & Relevant chapter section & Relevant literature & Suggested next diagnostic \\
\midrule
\endfirsthead
\caption[]{Literature-to-debugging crosswalk for the rebuilt DPF block,
continued.}\\
\toprule
Failure mode & Likely mathematical layer & Relevant chapter section & Relevant literature & Suggested next diagnostic \\
\midrule
\endhead
\bottomrule
\endlastfoot
ESS collapse or weight concentration & classical SMC degeneracy, proposal mismatch, or high-dimensional likelihood concentration & Chapter~\ref{ch:bf-particle-filters}; Chapter~\ref{ch:bf-dpf-particle-flow-foundations}, Sections~\ref{sec:bf-pff-motivation} and \ref{sec:bf-pff-stiffness} & \citet{gordon1993novel,doucet2001sequential,chopin2020introduction,bengtsson2008curse,snyder2008obstacles} & plot ESS by time and parameter region; compare bootstrap PF, flow proposal, and trusted low-dimensional reference; isolate whether collapse occurs before or after flow \\
Flow stiffness or ODE instability & homotopy transport and numerical integration of particle-flow ODE & Chapter~\ref{ch:bf-dpf-particle-flow-foundations}, Sections~\ref{sec:bf-pff-continuity}, \ref{sec:bf-pff-edh}, and \ref{sec:bf-pff-stiffness} & \citet{daumhuang2008,li2017particle,hu2021particle} & sweep flow step size and integrator tolerance; monitor endpoint error in the linear-Gaussian recovery case; record Jacobian spectra along $\lambda$ \\
EDH/LEDH mismatch & global Gaussian closure versus particle-local linearization & Chapter~\ref{ch:bf-dpf-particle-flow-foundations}, Sections~\ref{sec:bf-pff-edh}, \ref{sec:bf-pff-linear-gaussian-recovery}, and \ref{sec:bf-pff-ledh} & \citet{daumhuang2008,li2017particle,hu2021particle} & compare EDH and LEDH endpoint moments on the same particles; test whether local Jacobians change discontinuously across particles or parameter draws \\
PF-PF weight or likelihood disagreement & proposal density, change of variables, and correction weight & Chapter~\ref{ch:bf-dpf-pfpf}, Sections~\ref{sec:bf-pfpf-change-of-variables}, \ref{sec:bf-pfpf-corrected-weights}, and \ref{sec:bf-pfpf-what-restored} & \citet{li2017particle,doucet2001sequential,chopin2020introduction} & verify the pre-flow and post-flow proposal densities on a synthetic affine flow; compare corrected and uncorrected weights; check normalized log weights against a brute-force reference \\
Jacobian or log-det mismatch & determinant evolution under the flow map & Chapter~\ref{ch:bf-dpf-pfpf}, Section~\ref{sec:bf-pfpf-logdet} & \citet{li2017particle} and the change-of-variables argument in Chapter~\ref{ch:bf-dpf-pfpf} & compare integrated scalar log-det to finite-difference or automatic-differentiation Jacobian determinants in small dimension; check sign convention and time integration direction \\
Zero or misleading resampling gradients & hard categorical resampling or a changed soft-resampling law & Chapter~\ref{ch:bf-dpf-resampling-ot}, Sections~\ref{sec:bf-dr-discontinuous-map}, \ref{sec:bf-dr-soft-resampling}, and \ref{sec:bf-dr-bias-vs-diff} & \citet{zhumurphyjonschkowski2020,corenflos2021differentiable} & trace whether ancestor selection remains discrete; test nonlinear test-function bias under soft resampling; compare gradient path to the scalar actually evaluated \\
Sinkhorn marginal residuals or NaN gradients & entropy-regularized OT solve and finite-iteration numerical approximation & Chapter~\ref{ch:bf-dpf-resampling-ot}, Sections~\ref{sec:bf-dr-transport-view}, \ref{sec:bf-dr-entropic-ot}, and \ref{sec:bf-dr-debug-map} & \citet{villani2003topics,cuturi2013sinkhorn,peyre2019computational,schmitzer2019stabilized,corenflos2021differentiable} & monitor row and column residuals, transport-cost gap, $\varepsilon$ sensitivity, iteration budget, and log-domain stabilization; separate teacher-object error from solver error \\
Learned-OT permutation or extrapolation failure & learned set map approximating an OT teacher & Chapter~\ref{ch:bf-dpf-learned-ot}, Sections~\ref{sec:bf-learned-ot-teacher}, \ref{sec:bf-learned-ot-residuals}, and \ref{sec:bf-learned-ot-training-distribution} & \citet{corenflos2021differentiable,zaheer2017deep,lee2019set} & permute particle order and compare outputs; evaluate teacher-student residuals under sharp weights, larger $N$, shifted state dimension, and out-of-distribution parameter regions \\
DPF-HMC target mismatch & scalar target, gradient object, and posterior interpretation & Chapter~\ref{ch:bf-dpf-hmc-target-suitability}, Sections~\ref{sec:bf-dpf-hmc-contract}, \ref{sec:bf-dpf-hmc-rung-analysis}, and \ref{sec:bf-dpf-hmc-summary} & \citet{neal2011mcmc,andrieu2009pseudo,andrieu2010particle,betancourt2017conceptual} & audit whether $g(\theta)=\nabla_\theta\ell(\theta)$ for the same $\ell$; run finite-difference gradient checks; classify the rung as exact, corrected particle, relaxed, or learned surrogate before interpreting chains \\
Value object versus gradient object mismatch & implementation boundary between likelihood value, autodiff path, compiled execution, and sampler interface & Chapter~\ref{ch:bf-dpf-hmc-target-suitability}, Sections~\ref{sec:bf-dpf-hmc-contract}, \ref{sec:bf-dpf-hmc-recommendation}, and \ref{sec:bf-dpf-hmc-boundary} & \citet{neal2011mcmc,betancourt2017conceptual,greydanus2019hamiltonian} & compare value-only, autodiff, and finite-difference paths on fixed random seeds; test compiled versus eager execution; verify that accept/reject uses the same scalar whose gradient drove the trajectory \\
\end{longtable}
\end{center}

\section{Diagnostic order}
\label{sec:bf-dpf-debug-order}

The safest debugging order is from target definition outward:
\begin{enumerate}[label=(\arabic*)]
  \item define the scalar object being estimated or sampled;
  \item verify that the gradient corresponds to that scalar object;
  \item isolate whether errors enter before flow, during flow, in proposal
    correction, at resampling, in OT/Sinkhorn, or in learned transport;
  \item compare each layer with a special case whose answer is known;
  \item only then tune numerical tolerances or architecture.
\end{enumerate}
This order matters because many DPF failures can be made numerically smoother
without becoming mathematically correct.  A low-variance learned map, for
example, can still be the wrong target; a stable Sinkhorn loop can still
represent the wrong $\varepsilon$; and a clean HMC trace can still follow a
gradient that is not the derivative of the scalar being accepted.

\section{Coding-agent issue taxonomy}
\label{sec:bf-dpf-debug-agent-taxonomy}

For implementation work, the categories in
Table~\ref{tab:bf-dpf-debug-crosswalk} should be treated as routing labels:
\begin{itemize}
  \item \textbf{SMC degeneracy}: ESS, log-weight variance, proposal/reference
    comparison.
  \item \textbf{Flow transport}: endpoint recovery, stiffness, local Jacobian
    stability, discretization.
  \item \textbf{Proposal correction}: density-ratio algebra, Jacobian/log-det
    sign, corrected-weight normalization.
  \item \textbf{Differentiable resampling}: categorical discontinuity, relaxed
    law, nonlinear-functional bias.
  \item \textbf{OT/Sinkhorn}: marginal residuals, regularization, finite-iteration
    error, numerical stabilization.
  \item \textbf{Learned transport}: teacher residual, permutation equivariance,
    training-distribution coverage, posterior shift.
  \item \textbf{HMC target contract}: value-gradient consistency, target-status
    classification, compiled-path reproducibility.
\end{itemize}

These labels are also a guard against lane drift.  Debugging a learned-OT
failure should not trigger a rewrite of the particle-flow derivation unless the
diagnostic has shown that the teacher itself is wrong.  Debugging an HMC
acceptance failure should not trigger a resampling change until the
value-gradient contract has been checked.

\section{Boundary of this crosswalk}
\label{sec:bf-dpf-debug-boundary}

This chapter does not validate any DPF implementation.  It gives a structured
route from observed failures to mathematical suspects and source-backed
diagnostics.  The next work item after this enrichment round is therefore not a
generic rewrite of the chapters, but a set of targeted experiments that exercise
the rows of Table~\ref{tab:bf-dpf-debug-crosswalk}: linear-Gaussian recovery,
affine-flow Jacobian checks, Sinkhorn residual sweeps, learned-OT
out-of-distribution tests, and HMC value-gradient consistency tests.

\FloatBarrier
